\documentclass{report}
\usepackage{amsmath,amssymb}
\setlength{\parindent}{0mm}
\setlength{\parskip}{1em}
\begin{document}
\begin{center}
\rule{6in}{1pt} \
{\large Ruipeng Li \\
{\bf A Thick-Restart Lanczos algorithm with polynomial filtering for Hermitian eigenvalue problems}}

Center for Applied Scientific Computing \\ Lawrence Livermore National Laboratory \\ P O Box 808 \\ L-561 \\ Livermore \\ CA 94551
\\
{\tt li50@llnl.gov}\\
Yuanzhe Xi\\
Eugene Vecharynski\\
Chao Yang\\
Yousef Saad\end{center}

Polynomial filtering can provide a highly effective means of computing
all eigenvalues of a real symmetric (or complex Hermitian) matrix that
are located in a given interval, anywhere in the spectrum. This work
presents a technique for tackling this problem by combining a
Thick-Restart version of the Lanczos algorithm with deflation (`locking')
and a new type of polynomial filters. Thick restarting is employed to
limit the cost of orthogonalization. The polynomial filter that maps
eigenvalues within a given interval to eigenvalues with the largest
magnitude of the transformed problem, is obtained from a least-squares
approximation to an appropriately centered Dirac-$\delta$ distribution.
The resulting algorithm
can be utilized in a `spectrum-slicing' approach whereby a
very large number of eigenvalues and associated eigenvectors of the
matrix are computed by extracting eigenpairs located in different
sub-intervals independently from one another. Numerical experiments show
that such a construction yields effective polynomial filters which along
with a Thick-Restart Lanczos procedure enable desired eigenpairs to be
computed efficiently.


\end{document}
